\emph{(i).} Write $\tau = \tr h$ and $Q = Q(k)$. From $\mathcal{R}_{k}[h](B) = -\tfrac12\,k \wedge h(B \cdot k)$ (Theorem \claimref{synthesis.weyl-flat-variations-carry-only-the-stress-trace-and-the-whole-stress-clause-is-the-trace-free-channel}; for a plane-wave profile the curvature is this map times a function of $k \cdot x$) and $(e_{b} \wedge a) \cdot k = (a \cdot k)e_{b} - (e_{b} \cdot k)a$,
\[
\mathrm{Ric}(a) = \sum_{b} e^{b} \cdot \mathcal{R}_{k}[h](e_{b} \wedge a) = -\tfrac12\Bigl[(a \cdot k)\bigl(h(k) - \tau k\bigr) - Q\,h(a) + \bigl(k \cdot h(a)\bigr)k\Bigr],
\]
using $e^{b} \cdot (k \wedge v) = (e^{b} \cdot k)v - (e^{b} \cdot v)k$, and the scalar is $-(k \cdot h(k) - Q\tau)$. A displacement changes $h$ by $a \mapsto (a \cdot k)\xi + (a \cdot \xi)k$ and leaves the curvature unchanged (clause (i) of that theorem). With $\bar h = h - \tfrac12\,\id\,\tau$ it changes $\bar h(k)$ by $Q\xi$, so for $Q \neq 0$ there is exactly one displacement after which $\bar h(k) = 0$, that is $h(k) = \tfrac12\tau k$. Then $\mathrm{Ric}(a) = \tfrac12 Q\,h(a)$, the scalar is $\tfrac12 Q\tau$, and $\mathcal{G}_{k}[h] = \mathrm{Ric} - \tfrac12\,\id\cdot\tfrac12 Q\tau = \tfrac12 Q\,\bar h$. For a source $\mathcal{T}$ put $h = \mathcal{T} - \tfrac12\,\id\,\tr\mathcal{T}$: then $\bar h = \mathcal{T}$, $\bar h(k) = \mathcal{T}(k) = 0$, and $\mathcal{G}_{k}[h] = \tfrac12 Q\,\mathcal{T}$. If $\mathcal{G}_{k}[h] = 0$ with $Q \neq 0$, then after the displacement $\bar h = 0$, so $h = 0$: the kernel is the space of displacements, of dimension four since $\xi \mapsto \bigl(a \mapsto (a \cdot k)\xi + (a \cdot \xi)k\bigr)$ is injective for $k \neq 0$. So the frame equation $\mathcal{G} + 8\pi G\,\mathcal{T} = 0$ of Fact \claimref{synthesis.imported-einstein-cartan-matter-coupling}(a), linearised, has at each $k$ with $Q \neq 0$ the solution $h = \lambda(\mathcal{T} - \tfrac12\,\id\,\tr\mathcal{T})/Q$, unique up to a displacement, with $\lambda$ a constant fixed by $G$ and the sign conventions of the curvature and of $h$. The inverses of the wave operator (retarded, advanced, symmetric, time-ordered) differ only by solutions of the equation without source, and at $Q \neq 0$ there are none. At null $k$ the rank of $\mathcal{G}_{k}$ on the ten-dimensional space of symmetric maps is four (computed at two null momenta), so the kernel has dimension six: the displacements and two further directions, the transverse trace-free ones of clause (ii) of the cited theorem.

\emph{(ii).} By clause (v) of the cited theorem the first-order response of matter to a change of the field is the pairing of that change with its stress map, so the response of matter $1$ to the field of matter $2$ is a constant times $\ip{\mathcal{T}_{1}}{\mathcal{T}_{2} - \tfrac12\,\id\,\tr\mathcal{T}_{2}}/Q = B(\mathcal{T}_{1},\mathcal{T}_{2})/Q$ (displacements pair to zero with a source, clause (iii) there), symmetric in the two. The constant does not depend on the sources. For bodies at rest, $\mathcal{T}_{i}(a) = m_{i}(a \cdot e_{0})e_{0}$: $\ip{\mathcal{T}_{1}}{\mathcal{T}_{2}} = m_{1}m_{2}$, the traces are $m_{1}$, $m_{2}$, $B = \tfrac12 m_{1}m_{2}$, and $k = (0, \mathbf{k})$ has $Q = -|\mathbf{k}|^{2}$. The stored coupling of Proposition \claimref{field-theory.operator-newtonian-coupling} is $-Gm_{1}m_{2}/r$ between two bodies, with transform $-4\pi Gm_{1}m_{2}/|\mathbf{k}|^{2}$ (the transform of $1/r$ is $4\pi/|\mathbf{k}|^{2}$). Hence the constant is $8\pi G$. With the stress maps taken between states normalised to twice the energy per quantum (the convention of the cited cross-section, Theorem \claimref{synthesis.electron-positron-to-muon-pairs-in-the-spacetime-algebra-lsz-tree-cross-section-and-one-loop-unitarity}, $d\sigma/d\Omega = |\mathcal{M}|^{2}/(64\pi^{2}s)$), the interaction is the invariant amplitude; for bodies at rest that normalisation multiplies $B$ by $4m_{1}m_{2}$, and $\mathcal{M} = -4m_{1}m_{2}$ times the transform of the potential is the Born relation.

\emph{(iii).} The stress map of a massless scalar field is $(a \cdot \nabla\phi)\nabla\phi - \tfrac12 a\,Q(\nabla\phi)$ (clause (v) of the cited theorem), to which the curvature coupling adds $\xi\bigl(a\,\nabla^{2} - (a \cdot \nabla)\nabla\bigr)\phi^{2}$, the flat limit of the improved stress. Between the vacuum and two quanta each gradient gives the momentum times the field's unit, whose square is $-1$, both assignments of the two quanta contribute, and $\phi^{2}$ gives $2$ at total momentum $k$: the stated $\mathcal{T}$. Its trace is $-2p_{1} \cdot p_{2} + 4p_{1} \cdot p_{2} + 2\xi(Q(k) - 4Q(k)) = (1 - 6\xi)s$, with $2p_{1} \cdot p_{2} = s$. For the vector field the stress is $-\tfrac12 FaF$ (Fact \claimref{synthesis.imported-maxwell-stress-spacetime-algebra}) and the two-quantum part of $F = F_{1} + F_{2}$ is the stated cross term. Two forms used here are not in the cited clause, which covers minimally coupled matter only, and are standard, cited from memory: the improvement term of the scalar stress just used, and the symmetric stress of the Dirac field (the form with one derivative acting to either side, symmetrised with the current's direction; it is the form the published calculations use). For the fermion that stress has one derivative and one current: between the vacuum and the pair the derivatives give $d = p_{1} - p_{2}$ and the current its transition value $j$, with $j \cdot p_{1} = j \cdot p_{2} = 0$ on shell; the spin sum of $(j \cdot w)^{2}$ for one helicity is half the trace of the cited theorem's clause (1) at zero mass, $\tfrac12 \cdot 4\br{p_{2}wp_{1}w}_{0}$.

Scalars to scalars, $\xi = 0$. With $A_{1}(a) = (a \cdot p_{2})p_{1} + (a \cdot p_{1})p_{2}$ and $A_{2}$ likewise from $p_{3}, p_{4}$: $\ip{\mathcal{T}_{1}}{\mathcal{T}_{2}} = \ip{A_{1}}{A_{2}} = 2\bigl[(p_{1} \cdot p_{3})(p_{2} \cdot p_{4}) + (p_{1} \cdot p_{4})(p_{2} \cdot p_{3})\bigr] = \tfrac12(t^{2} + u^{2})$, the terms with $\id$ cancelling ($\tr A_{i} = s$), and the traces are $s$, so $B = \tfrac12(t^{2} + u^{2} - s^{2}) = -tu$. In the centre-of-mass frame the spatial part of $\mathcal{T}_{1}$ is $2\mathbf{p}\mathbf{p}$ minus a multiple of the identity; its trace-free part pairs with that of $\mathcal{T}_{2}$ to $\tfrac{s^{2}}{4}(\cos^{2}\theta - \tfrac13)$, and the rest is $-\tfrac16$ of the product of the traces, which the term in $\xi$ changes to $(1 - 6\xi_{1})(1 - 6\xi_{2})s^{2}$ without touching the trace-free part ($(a \cdot k)k - Q(k)a$ is a multiple of the spatial identity in that frame).

Fermion pair to scalars. $\ip{\mathcal{T}_{f}}{\mathcal{T}_{2} - \tfrac12\id\tr\mathcal{T}_{2}} = \tfrac12\,j \cdot w$ with $w = \mathcal{T}_{2}(d) - \tfrac12 s\,d$; since $j \cdot d = j \cdot k = 0$ only the part $-(d \cdot p_{4})p_{3} - (d \cdot p_{3})p_{4}$ of $w$ counts, and $d \cdot p_{4} = \tfrac12(t - u) = -d \cdot p_{3}$, $j \cdot p_{4} = -j \cdot p_{3}$, so $B = -\tfrac12(t - u)\,j \cdot p_{3}$ and $\sum B^{2} = \tfrac14(t - u)^{2}\cdot 2\br{p_{2}p_{3}p_{1}p_{3}}_{0} = \tfrac14(t - u)^{2}\cdot 4(p_{1} \cdot p_{3})(p_{2} \cdot p_{3}) = \tfrac14(t - u)^{2}tu$.

Vectors to scalars. With transverse real polarisations, $F_{1}aF_{2} + F_{2}aF_{1}$ pairs with $A_{2}$ to $-tu\,(\epsilon_{1} \cdot \epsilon_{2}) - 2s(\epsilon_{1} \cdot p_{3})(\epsilon_{2} \cdot p_{3})$ up to the overall factor; summing the squares over the two polarisations of each quantum gives $2t^{2}u^{2}$ (computed exactly at four configurations).

\emph{(iv).} In the centre-of-mass frame at fixed $s$ each source is a fixed symmetric map built from its pair's axis, and the bracket of the two is a polynomial of degree at most two in the cosine and sine of the angle $\theta$ between the axes; the sum of squares over polarisations is even in the sine by reflection in the plane of the two axes, hence a polynomial of degree at most four in $\cos\theta$. It is computed exactly at seven rational angles in each of the nine channels (and for both relative helicities of fermions to fermions), the coefficients above degree four vanish, and the integral over the sphere is taken exactly. For the three channels into scalars this agrees with (iii): $\int t^{2}u^{2}\,d\Omega = \tfrac{2\pi}{15}s^{4}$, $\int\tfrac14(t - u)^{2}tu\,d\Omega = \tfrac{\pi}{30}s^{4}$, and the factors $\tfrac14$, $\tfrac12$, $\tfrac12$ for identical quanta give $\tfrac{\pi}{30} = \tfrac{\pi}{180}(1 + 5)$, $\tfrac{\pi}{60}$, $\tfrac{\pi}{15}$. The product form reflects that away from the trace part only total angular momentum two is exchanged; the weights $1 : 3 : 12$ are those of the three fields in the spin-two part of the two-point function of the stress (standard; cited from memory, not used).

\emph{(v).} The rate per volume of a transition $1 + 2 \to 3 + 4$ from a bath with occupation functions $f_{1}, f_{2}$, the final species being absent, is $\int d\Pi_{1}\,d\Pi_{2}\,f_{1}f_{2}\int|\mathcal{M}|^{2}\,d\Phi_{2}$ with $d\Pi = d^{3}p/((2\pi)^{3}2E)$ and, for massless quanta, $d\Phi_{2} = d\Omega/(32\pi^{2})$, a factor $\tfrac12$ for each pair of identical quanta (the standard collision integral, in the normalisation of the cited cross-section). With $d\Pi_{1}\,d\Pi_{2} = E_{1}E_{2}\,dE_{1}\,dE_{2}\,d\cos\theta_{12}/(32\pi^{4})$ after the trivial angles, $|\mathcal{M}|^{2} = (8\pi G)^{2}B^{2}/s^{2}$ and (iv), the inner integral is $W_{ij}s^{2}/\bar{M}_{P}^{4}$; $s = 2E_{1}E_{2}(1 - \cos\theta_{12})$ gives $\int s^{2}\,d\cos\theta_{12} = \tfrac{32}{3}E_{1}^{2}E_{2}^{2}$; and $\int_{0}^{\infty}E^{3}(e^{E/T} \mp 1)^{-1}dE = \tfrac{\pi^{4}}{15}T^{4}$, $\tfrac{7\pi^{4}}{120}T^{4}$ for Bose and Fermi occupation (from $\zeta(4) = \pi^{4}/90$). So
\[
\Gamma_{j} = \frac{1}{1024\pi^{6}}\cdot\frac{32}{3}\cdot\frac{T^{8}}{\bar{M}_{P}^{4}}\Bigl[\frac{\pi^{8}}{225}\sum_{\text{bosons}}W_{ij} + \frac{49\pi^{8}}{14400}\sum_{\text{fermions}}W_{ij}\Bigr],
\]
which is the stated $\alpha_{j}$ on inserting $W_{ij}$. The numbers follow: for a scalar, $\bigl[(148 + 20)/225 + 49 \cdot 45/4800\bigr]/17280 = 5789/82944000$; for a chiral fermion $\bigl[3 \cdot 148/225 + 49 \cdot 135/4800\bigr]/17280 = 16087/82944000$; a real vector is four times that; the total over four scalars, $45$ chiral fermions and twelve vectors is $1519247/82944000$. At $s = 0$ the integrand $B^{2}/s^{2}$ is $s^{2}$ times a bounded function of the angles, so the prescription at $Q(k) = 0$, where alone the inverses of (i) differ, does not enter. Hypotheses of the formula, besides (OS): tree level; a free bath whose quanta are massless at the temperature in question (for the content $(4, 45, 12)$, above the electroweak scale, two polarisations per vector), with no thermal masses; zero chemical potentials; no occupation of the produced species, so no final-state factor. The effective interaction of second order built from the retarded field is the principal-value kernel, which is the time-ordered one away from $Q(k) = 0$.

\emph{(vi).} Arithmetic: with $N_{w}$ replaced by $2N_{w}$ the scalar bracket is $(168/225 + 49 \cdot 90/4800)/17280 = 3997/41472000$ and, for $c_{j} = 6$, $\tau_{j} = 0$, $(6 \cdot 148/225 + 49 \cdot 6 \cdot 90/4800)/17280 = 11351/20736000$. These are the two coefficients printed in arXiv:2112.15214, whose squared amplitudes for scalar production, $\tfrac14 t^{2}u^{2}/s^{2}$, $\tfrac14\,tu\,(s + 2t)^{2}/s^{2}$, $\tfrac12 t^{2}u^{2}/s^{2}$ in units of $\bar{M}_{P}^{-4}$ with the symmetry factors included (its equations (83) to (85), where $tu$ is written $-t(s+t)$), are those of (iii) with the fermion's sum taken over both helicities; with multiplicity $45$ that weights each chiral fermion as a Dirac fermion. $663552000 = 16 \cdot 41472000$. The origin of the sixteen is in the earlier paper itself, read on the web in this run's review: its coupling is $h_{\mu\nu}T^{\mu\nu}/(2\bar{M}_{P})$ (its equation (3)), its propagator is that of a canonically normalised field (appendix, equation (17)), and its squared amplitudes (appendix, equations (25) to (27)), $\tfrac1{16}t^{2}u^{2}/s^{2}$, $\tfrac1{32}tu(t-u)^{2}/s^{2}$, $\tfrac18 t^{2}u^{2}/s^{2}$ without symmetry factors and with the fermion summed over both helicities, are one sixteenth of the canonically normalised ones; inserted in its rate integral with the factors $\tfrac14$, $\tfrac12$, $\tfrac14$ for identical quanta they give its equation (31), the stored value. Half the canonical coupling at each of two vertices is a factor four in the amplitude and sixteen in the rate. The later paper prints its value with a citation of the earlier one and no remark on the difference.

\emph{(vii).} The expectation value of the stress of a homogeneous isotropic bath is homogeneous and isotropic, so the field it sources is the flat field rescaled by a function of time, with vanishing Weyl curvature. By Fact \claimref{synthesis.imported-conformal-covariance-of-massless-fields-and-local-conformal-flatness}(b) the solution spaces and the algebras of the chiral fermion, the vector field and the conformally coupled scalar are those of the flat field, and the state that is the vacuum before stays the vacuum in the rescaled modes: no pairs. For a scalar with $\xi \neq \tfrac16$ the rescaled mode equation acquires the term $(\xi - \tfrac16)$ times the scalar curvature, which vanishes when the source has trace-free stress, as a bath of massless quanta has. That the conformal vacuum stays the vacuum is the standard argument (Parker; cited from memory) and is not in the cited Fact, which supplies the identification of the solution spaces and algebras. A real bath is not exactly trace-free (masses, running couplings, the trace anomaly), so the statement holds up to terms suppressed by the trace of the bath's stress and by the ratio of the expansion rate to the temperature. Fluctuations of the stress about its mean are what (OS) keeps and the mean-field source drops.

\emph{Reading.} The factors quoted: for the mirror, $0.568$ against the stored bracket $0.053$ to $0.078$, ratios $10.7$ and $7.3$, fourth roots $1.81$ and $1.64$; for nine massless vectors, $9 \times 0.02405 = 0.2165$ against $0.021$, ratio $10.3$, fourth root $1.79$.
